\section{Introduction}\label{sec:introduction}

We fix the lattice spacing and use unit lattice coordinates. The photon
field is a real one-cochain on $\Z^3$. Write $d$ for the lattice
coboundary, orient a plaquette by the ordered pair $(e_j,e_k)$, $j<k$,
and let $\mathcal P_2$ contain the plaquettes with these orientations.
The space in which the volume limit is taken is
\[
 \Omega=\{F\in\R^{\mathcal P_2}:dF=0\},
\]
with the product topology. Reversing a plaquette's orientation negates
its field value. In particular, pullback by a lattice isometry acts on
$F$ as a two-cochain, including the orientation sign.

Here is the finite model. Let
$W=\Z^3\rtimes\bigl(O(3)\cap\operatorname{GL}(3,\Z)\bigr)$ and
$C_0=\{0<x_3<x_2<x_1<1/2\}$. Fix $v_0\in C_0$, and let
$s_0,s_1,s_2,s_3$ be the reflections in its walls
$x_3=0$, $x_2=x_3$, $x_1=x_2$, $x_1=1/2$, respectively.
The fermion sites are $\mathcal X=Wv_0$, and the four neighbours
of $x=gv_0$ are $x^{(i)}=gs_iv_0$. Write $\pi(x)$ for the nearest
point of $\Z^3$ to $x$. With $\gamma_j=\operatorname{diag}(\sigma_j,-\sigma_j)$,
where $\sigma_j$ are the Pauli matrices, put $\gamma_v=\sum_jv_j\gamma_j$.
For charge $q$, Wilson parameter $\rho$, weights $w_i$, and mass
parameter $M$, the chamber operator is
\[
 (\mathcal D_A\psi)(x)
 =M\psi(x)-\sum_{i=0}^3\frac{w_i}{2}
   \left(\rho-\gamma_{(x^{(i)}-x)/|x^{(i)}-x|}\right)
   e^{iqA(\pi(x),\pi(x^{(i)}))}\psi(x^{(i)}),
\]
where $A(y,x)=-A(x,y)$ and $A(x,x)=0$.

On the cubic torus $(\Z/n\Z)^3$, the photons are periodic and
$\mathcal D_A^{C(n)}$ acts on spinors satisfying
$\psi(x+ne_j)=-\psi(x)$ for $j=1,2,3$. Put $r_q=2\pi/q$.
For a bounded Borel function $G$ on $\Omega$, define
\[
 \mu_n(G)=\frac1{Z_{C(n)}[1]}
 \sum_{\substack{z\in C^2((\Z/n\Z)^3;\Z)\\dz=0}}
 \int_{[0,r_q)^{3n^3}}
 G(\widetilde F)\,e^{-\|F\|^2/2}
       \det\mathcal D_A^{C(n)}\,\dd A,
 \qquad F=dA+r_qz.
\]
Here $\widetilde F$ is the periodic extension of $F$, the norm sums
over the torus plaquettes, and $Z_{C(n)}[1]$ is the same expression
without $G$. Thus every closed integer cochain is summed, with no
restriction on its flux. The determinant is nonnegative and
$0<Z_{C(n)}[1]<\infty$; these facts are proved in
\cref{prop:reg-determinant-sign,prop:reg-partition}.

\begin{theorem}[Reflection positivity in infinite volume]\label{thm:main}\label{thm:cubic-limits}\label{thm:diag-lattice-limits}
Fix $q>0$, $0<\rho<1$, $M\in\R$ and $w_0,\ldots,w_3>0$.
The family $(\mu_n)_{n\ge4}$ is tight on $\Omega$. Let
$n_i\ge4$ be any sequence of integers tending to infinity, and let
$\nu$ be any subsequential weak limit of $(\mu_{n_i})$. Then:
\begin{enumerate}
\item[\textup{(i)}] For every finitely supported real two-cochain $\beta$,
\[
 \int_\Omega e^{\langle\beta,F\rangle}\,\dd\nu(F)
       \le e^{\|\beta\|^2/2},
 \qquad
 \langle\beta,F\rangle=\sum_{\sigma\in\mathcal P_2}\beta_\sigma F_\sigma,
 \qquad \|\beta\|^2=\sum_{\sigma\in\mathcal P_2}\beta_\sigma^2.
\]
\item[\textup{(ii)}] The law $\nu$ is invariant under lattice translations,
signed coordinate permutations, and $F\mapsto-F$.
\item[\textup{(iii)}] The law $\nu$ is invariant and reflection positive
at every mirror
\[
 x_j=c\quad(c\in\tfrac12\Z),\qquad
 x_j\pm x_k=c\quad(c\in\Z,\ j<k).
\]
More precisely, write a mirror as $\Pi=\{\ell(x)=c\}$ and let
\[
 \theta_\Pi x
 =x-\frac{2(\ell(x)-c)}{|\nabla\ell|^2}\nabla\ell.
\]
Choose either open side $H$ of $\Pi$, and let $\Sigma_H$ be the
sigma-algebra generated by the values $F_\sigma$ for which every
corner of $\sigma$ lies in $H$. Then
\[
 \int_\Omega\overline{A(\theta_\Pi^*F)}A(F)\,\dd\nu(F)\ge0
 \qquad\bigl(A\in L^2(\Sigma_H,\nu)\bigr).
\]
\end{enumerate}
\end{theorem}

There are three coordinate and six diagonal mirror families in
\textup{(iii)}. The support convention excludes plaquettes on or
crossing the mirror. For example, reflection in $x_1=0$ changes the
signs of $F_{12}$ and $F_{13}$ and preserves the sign of $F_{23}$.
These signs occur both in invariance of the law and in its reflected
form.

\subsection{The comparison argument}

The coordinate mirrors are compatible with every cubic torus. Their
finite-volume reflected forms, together with a Gaussian bound,
give tightness and coordinate reflection positivity of every weak
limit. A diagonal mirror requires a different period lattice.
The task is therefore to obtain diagonal positivity for the limit
of the original cubic laws.

Reflected slab operators give a comparison with a diagonal torus
of the same volume. Insert a real local trigonometric polynomial
$O$ in one cell, reflect it into all cells, and tilt the diagonal
law by the exponential of $t$ times their sum. The resulting law
has reflection positivity at the cell walls. Its expectation of
$O$ bounds the cubic expectation, up to a pressure error divided
by $t$. The pressure error tends to zero as the volume grows, so
the volume limit must precede $t\downarrow0$.

The perturbation occupies the entire torus. Its pointwise bound
does not control local expectations uniformly in volume. We obtain
that control by amplifying local insertions in two cell structures
and applying a coordinate operator bound. Heat regularization then
proves approximate reflection symmetry on two adjacent fermion
layers. Fermion elimination and an imaginary contour shift reduce
each curvature character to these two layers, with an error of
order $t$ independent of the volume. After the ordered limits,
the comparison law is invariant under coordinate link reflections
and hence under translations by $2\Z^3$.

These translations allow a given observable to be placed inside a
cell with any prescribed diagonal wall. For each real local
trigonometric polynomial $O$ we obtain a comparison law $\nu_O$
with positivity at that wall and
$\nu(O)\le\nu_O(O)$. The comparison law may depend on $O$.
The laws satisfying the Gaussian bound and positivity at the wall
form a compact convex set, and separation places $\nu$ itself in
that set.

The repeated reflections require divisible side lengths, whereas
the theorem allows arbitrary integers $n_i$. An unequal-height
slab estimate compares each original cube to a nearby divisible
companion, whose side $m_i$ satisfies $m_i/n_i\to1$. The H\"older-weighted
volume of the exceptional slabs is a vanishing fraction of the original
volume, and a uniform pressure bound controls their contribution. All auxiliary states
are constructed on the companions; the final separation argument
still applies to the prescribed limit of the original cubes.


This paper gives a self-contained reorganization of the fixed-lattice
arguments in \cite[Sections~3--11]{ChamberQED}. The volume-limit
theorem, including arbitrary cubic side lengths, is already stated
there in Theorem~1.2 and Corollary~11.19. We reproduce the finite-volume
identities and limiting arguments, retain the unequal-slab comparison
needed for arbitrary side lengths, and spell out the extension of the
reflection form to the full open-side $L^2$ space. Repeated reflection and the
chessboard method belong to the general theory of Fr\"ohlich, Israel,
Lieb, and Simon~\cite[Section~4]{FILS1978}. Here the slab trace must also
retain the photon flux sum and the fermion boundary character when
the period lattice changes. The local operator and layer-elimination
arguments establish the additional symmetry needed for that comparison.
The theorem concerns fixed lattice parameters and subsequential volume
limits; it asserts neither uniqueness nor a continuum limit.
